%% =====================================================================
%%  T-17-06  The Slow Reform
%%  -------------------------------------------------------------------
%%  One idea per file. Core is mandatory. Slots in canonical order.
%%  Max 700 words. No layout commands. No filler. New source material
%%  for this entry goes in sources/B3/T-17-06.tex -- never by
%%  rewriting this file (docs/RULES.md R0.1).
%%  =====================================================================
%% @kind: topic
%% @id: T-17-06
%% @title: The Slow Reform
%% @subtitle: Preach change, but reform slowly --- innovation is trauma in a suit
%% @chapter: c17
%% @order: 60
%% @status: stable
%% @sources: B3
%% @updated: 2026-09-19

\topic{T-17-06}{The Slow Reform}{Preach change, but reform slowly --- innovation is trauma in a suit}

\begin{core}
Everyone grants the need for change in the abstract; at the scale of the actual Tuesday, people are habit's subjects, and too much innovation is trauma that leads to revolt. The delivery discipline for anyone new to a position: honour the house's way of doing things, visibly, before touching it, and let change arrive dressed as a gentle improvement on the past --- the reform that lands is the reform nobody had to survive.
\end{core}
\begin{mechanism}
B3's law runs on habit's protective function: routines are how people bank identity and predictability, and an outsider's reform is experienced as the theft of both. The transgression case is Cardinal Wolsey --- the minister who opposed the king's change was destroyed by it --- but the working cases all follow the same pattern: the successful reformer appeals to precedent, frames the new as restoration (back to the founders' way, back to what always worked), and sequences change so no single season contains enough novelty to organise against. The mechanism is reactance management: threatened habit triggers defence of the habit regardless of merits, so the skilful mover never triggers the threat --- each step is small enough to be adjudicated on its own merits, and the accumulated direction is only visible on a timescale long enough that opposition lacks a rally point. The source's corollary for newcomers: respect the old way visibly first; the audience must see you honouring the house before they will let you move the furniture.
\end{mechanism}
\begin{conditions}
Strongest for outsiders entering established systems --- new executives, new partners, reforms of institutions, in-laws --- and strongest of all where the audience's identity is fused with the practices being changed. It weakens in crisis, where habit is already broken and wholesale change is the only sale available, and it fails entirely against systems whose practices are actively harmful --- the slow reform of a burning house is arson by pacing.
\end{conditions}
\begin{application}
  \item On arrival, change nothing visible for a season. Learn the rituals, and be seen keeping them.
  \item Frame every change as restoration: find the precedent, quote the founder, name the tradition your reform fulfils.
  \item Sequence: one change per season, each small enough to be judged alone, none big enough to organise against.
  \item Let the house's own people propose what you need done. Authorship is the shield; supply the draft, accept the credit's loss.
  \item Keep a private map of the destination. Slow reform without a destination is drift with good manners.
\end{application}
\begin{feedback}
  \item Working: each change is accepted as obvious within a month. The pacing is right.
  \item Working: the old guard repeats your framing to newcomers. Restoration language has been adopted.
  \item Failing: a coalition forms around what you have not yet touched. The novelty per season crossed the revolt threshold.
  \item Failing: you have changed nothing in a year and the house is proud of you. The map and the manners have traded places.
\end{feedback}
\begin{failure}
The method's price is paid twice: by the reformer, who must let others author their work and be patient past their own enthusiasm; and sometimes by the house, where the slow frame is used to bury a change the members would have chosen fast if asked. The slow reform is also hostage to your tenure --- leave before the sequence completes and the house reverts, having learned only that outsiders flatter.
\end{failure}
\begin{countermeasures}
  \item When a newcomer praises your traditions effusively, watch what they change in month six. The respect is often scaffolding (T-13-03).
  \item Notice restoration framing in politics and management alike: ask what the golden age being quoted actually was. The past is being edited to sell the present.
  \item In your own house, separate pacing from direction: the slow pace is legitimate; the invisible destination is not. Demand the map.
  \item If you are the impatient reformer, remember the law's arithmetic: one revolt costs more years than ten delayed quarters.
\end{countermeasures}

\sources{B3}
\seesources
\seealso{T-12-02, T-17-04, T-21-05}
